\chapter{Para Além da Crença}

\lettrine{A}{\space partir do aparecimento} dos cinco \emph{khandhas} -- \emph{rūpa, vedanā, saññā,
  saṅkhārā, viññāṇa} -- e a crença inquestionável de que estes são o próprio eu
-- sempre parece que a mente está no corpo, certo? Para a maioria das pessoas, se
dissermos ``Onde está a sua mente?'', elas apontarão para a cabeça ou para o
coração.

Mas se investigarmos como as coisas são, seguindo os ensinamentos do Buddha,
começaremos a perceber que o corpo está na mente. A mente é o que vem primeiro
realmente -- o corpo é apenas o receptor. É um receptor sensível, como um rádio
ou radar. Não é uma pessoa, não é nada mais que um instrumento.

Quando essa visão de vivenciar-se dentro dos cinco \emph{khandhas} é percebida e
abandonada, então há uma percepção do que podemos chamar de `realidade imortal',
da imortalidade. Essas palavras subentendem `além do condicionado' e a
impossível capacidade de conceber a realidade imortal, não é? Podemos apontar
para uma palavra como `imortalidade' ou `imortal' ou `incondicionado', mas para
lá disso não há mais nada que se possa dizer sobre isso porque as palavras são
elas próprias condicionadas e mortais. Palavras, conceitos, percepções,
concepções, só são adequadas ao mundo condicionado. Enquanto estivermos apegados
a pensamentos e conceitos, a pontos de vista e opiniões, não importa quão
inteligentes e altruístas sejam esses pontos de vista, esse mesmo apego nos
ligará ao reino condicionado -- continuaremos a renascer nele. Continuamos à
procura do incondicionado no condicionado, continuamos à procura de Deus na
condição mortal, na natureza mutável da consciência sensorial, apenas para nos
sentirmos totalmente frustrados e desiludidos. Então inclinamo-nos a suportar
essa `visão da alma' por um tipo de crença teimosa.

As crenças não mudam, ou mudam? Podemos acreditar exactamente nas mesmas coisas
em que acreditávamos com cinco anos, agora com cinquenta! Essa crença é o apego
a uma percepção.

Existem algumas crenças que são muito boas, bonitas e sentimentais. A visão
romântica e sentimental da vida oferece uma bela imagem em que ainda podemos
acreditar mesmo aos 80 anos! Quando a minha avó morreu aos 75 anos, ela ainda
tinha o desenvolvimento emocional de uma moça de 16 anos. Quando ela morreu,
tinha um namorado chamado `Hercules Cavalier' que era seu gigolô! Ela ainda
tinha o mesmo tipo de desejo romântico de quando era uma adolescente de 16 anos!
Então, mesmo que aos 75 anos fosse um desastre físico, ainda assim a mente ainda
estava ligada a essas belas imagens de juventude.

Nós assumimos, acreditamos e nunca questionamos o preconceito e as ideias fixas
a que estamos agarrados, que nunca mudam. Elas não mudam; continuamos a
reafirmar repetidamente as mesmas velhas coisas. E é por isso que surgem tantos
problemas políticos porque muitas pessoas apegam-se a ideias políticas em vez de
tentarem estar cientes das necessidades de determinado tempo e lugar.

Quanta violência, malícia e maldade são perpetradas apenas em nome da
propriedade! E extremas: ``essa é minha terra, saia da minha terra''. Vemos isso
constantemente, nos intermináveis problemas de fronteira dos países. Depois
vê-se a mesquinhez do coração, o não querer deixar as pessoas entrarem ou não
querer deixar as pessoas saírem! Por causa da crença inquestionável -- ``esta é
a minha casa, a minha família, a minha esposa, o meu marido, os meus filhos,
meu, meu, meu\ldots{}''

Ao longo de vinte anos de meditação, posso ver que muitos apegos, obsessões e
tendências desapareceram ao permitir que as coisas cessem. O processo tem sido o
de deixar as coisas irem em vez de acreditar, agarrar e renascer em padrões de
pensamento e desejos sem fim. Quando vemos a vida como apenas uma passagem,
então não nos seguramos a ela. Não nos tornamos mesquinhos e egoístas porque
percebemos que nada vale a pena guardar -- uma riqueza material, propriedade,
status ou valores mundanos, ou qualquer coisa. Não há nada que valha a pena nos
chatearmos tanto, porque não nos pertence na realidade seja de que maneira for.
Claro que podemos acreditar que as coisas são nossas. Mas investigando na
realidade, ao observarmos a forma como a mente realmente é, vemos que nada na
verdade nos pertence de qualquer forma -- não há ninguém para possuir nada!

Agora, reflectindo sobre o corpo na mente, esse apego muda. Temos de começar a
contemplar o que a mente realmente é -- porque o nosso corpo na realidade não
está no nosso cérebro, certo? É ridículo! Nem o nosso coração que bombeia o
sangue. Eles estão no nosso corpo.

Eu estava de pé esta noite e olhava para o crepúsculo, para as árvores, as
árvores despidas nas extremas de Amaravati, simplesmente contemplando que as
árvores estão na mente e que as árvores são conscientes.

Há um certo nível de consciência em toda a vida, no facto de haver receptividade
ao meio ambiente; e as árvores são muito receptivas ao ambiente em que se
encontram. Começamos a mudar a percepção da mente para uma consciência que
permeia tudo. Então não é apenas uma mente humana, há algo mais para além disso.
Mas no budismo isso nunca é nomeado, nunca tentamos formar um conceito sobre
isso. Em vez disso, contemplamos a totalidade, toda a sensibilidade, o reino
sensorial e o que realmente envolve. E isso temos de contemplar a partir da
nossa própria capacidade de sermos conscientes e sentirmos, mas não em termos de
``eu'' e ``meu'' -- ``eu sinto essas coisas, mas ninguém mais sente'', ou
``somente os seres humanos conseguem e os animais não'', ou ``Só os mamíferos o
conseguem, não os répteis", ou "Só os reinos animal e dos insectos conseguem
duvidar, não o das plantas". A consciência não implica pensamento, mas implica
receptividade / sensibilidade a todo o sinal que vem do exterior, a tudo o que
lhe chega. Começamos a ver que a consciência é um sistema universal, vital e
mutável; é como uma reunião plenária, plena de todas as possibilidades, de todos
os potenciais de forma, do que pode ser criado. Tudo aquilo que conseguimos
pensar, podemos ver isso em termos da capacidade humana de imaginar, através da
qual podemos criar todos os tipos de fantasias que ganham forma material.

Mas o maior, mais profundo e significativo potencial humano é negligenciado pela
maioria das pessoas, e este é a capacidade de entender a verdade da realidade
como é, ver o Dhamma, estar livre de todas as ilusões.

Quando estiverem a contemplar a realidade, comecem a reflectir sobre onde não há
eu. Sempre que há a cessação do eu, há apenas clareza, conhecimento e
contentamento -- sentimo-nos à vontade e em equilíbrio. Leva um certo tempo até
sermos capazes de desistir de todos os esforços e tendências inquietas do corpo
e da mente. Mas, passado um tempo, isso cessará e há uma clareza real, uma paz
satisfeita. E aí também não há ego, nem ``eu'' e ``meu''. Podemos contemplar
isso.

Devemos compreender que temos que aprender através de uma forma totalmente
humilde, sem cobiçar sucesso em nada do que estamos a fazer nesta meditação,
nunca sendo bem sucedidos, nunca conseguindo o que queremos; se conseguirmos o
que queremos, perdemos imediatamente. Temos que ser totalmente humildes na forma
em que a nossa visão egocêntrica é abandonada voluntariamente, graciosamente,
humildemente. É por isso que, na meditação, quanto mais esta vem de uma força de
vontade baseada numa visão de si mesmo como ``eu vou alcançar e vou obter'', é
claro que só podemos esperar fracasso e desespero, porque a meditação não é uma
busca mundana. Em situações mundanas, se somos inteligentes e fortes, talentosos
e com oportunidades, existindo condições, podemos abrir caminho e termos grande
sucesso, certo?
